\section*{Bab \ref{ch:b2:mantle-melting} --- Pelelehan Mantel}

\begin{solution}{exo:b2:mantle-melting:1}
(a) Olivin lebih dari \qty{90}{\percent}: \omterm{def:b2:mantle-melting:peridotites}{dunit}. (b) Olivin dan ortopiroksen
tanpa klinopiroksen: \omterm{def:b2:mantle-melting:peridotites}{harzburgit}. (c) Olivin lebih dari \qty{40}{\percent} dengan
kedua piroksen: \omterm{def:b2:mantle-melting:peridotites}{lherzolit}. (d) Olivin kurang dari \qty{40}{\percent}: piroksenit
(websterit olivin, dengan kedua piroksen). (e) $57/98 = \qty{58.2}{\percent}$,
$28/98 = \qty{28.6}{\percent}$, $13/98 = \qty{13.3}{\percent}$: \omterm{def:b2:mantle-melting:peridotites}{lherzolit}.
\end{solution}

\begin{solution}{exo:b2:mantle-melting:2}
(a) Likuidus pada \qty{60}{\percent} anortit adalah \qty{1390}{\celsius}, di dalam
medan anortit: anortit mengkristal pertama. (b) Diopsid muncul pada
eutektik, \qty{1270}{\celsius}, ketika cairan mengandung \qty{42}{\percent} anortit.
(c) Cairan pertama dari campuran apa pun kedua mineral itu adalah cairan eutektik,
\qty{42}{\percent} anortit, pada \qty{1270}{\celsius}.
\end{solution}

\begin{solution}{exo:b2:mantle-melting:3}
(a) $\num{4.0e-5}\times1623/(3300\times1000) = \qty{1.97e-8}{K/Pa} =
\qty{19.7}{K/GPa}$. (b) $\num{e9}/(3300\times9.81) = \qty{3.09e4}{m}$, atau
\qty{30.9}{km} per gigapascal; $19.7/30.9 = \qty{0.64}{K/km}$, sekitar lima puluh kali
lebih kecil daripada \qty{30}{K/km} di dekat permukaan, tempat kalor dihantarkan secara konduksi, bukan
dibawa oleh aliran.
\end{solution}

\begin{solution}{exo:b2:mantle-melting:4}
(a) $T_{\mathrm s}(0) = \qty{1085.7}{\celsius}$; $T_{\mathrm s}(1) = 1085.7 +
132.9 - 5.1 = \qty{1213.5}{\celsius}$; $T_{\mathrm s}(2) = 1085.7 + 265.8 - 20.4 =
\qty{1331.1}{\celsius}$; $T_{\mathrm s}(3) = 1085.7 + 398.7 - 45.9 =
\qty{1438.5}{\celsius}$. (b) $132.9 - 2\times5.1\times2 = \qty{112.5}{K/GPa}$.
(c) Pada \qty{3}{GPa}, \qty{1400}{\celsius} berada di bawah solidus,
\qty{1438.5}{\celsius}; ketika naik ke \qty{2}{GPa} menyusuri adiabat, mantel hanya mendingin
sekitar \qty{20}{\celsius}, menjadi kira-kira \qty{1380}{\celsius}, sedangkan solidus turun menjadi
\qty{1331.1}{\celsius}: batuan kini berada di atas solidusnya dan sebagiannya meleleh.
\end{solution}

\begin{solution}{exo:b2:mantle-melting:5}
(a) $1350 + (1277.5 - 1350)\times(30 - 20)/(40 - 20) = \qty{1314}{\celsius}$.
(b) Komposisinya terletak di sisi diopsid dari eutektik: diopsid.
(c) Cairan mencapai eutektik dengan \qty{42}{\percent} anortit; anortit
tetap berada dalam cairan, sehingga fraksi cairannya $30/42 = 0.714$ dan $1 - 0.714 =
\qty{29}{\percent}$ sistem telah mengkristal sebagai diopsid. (d) Semuanya
mengkristal: \qty{70}{\percent} diopsid dan \qty{30}{\percent} anortit, yaitu
komposisi totalnya.
\end{solution}

\begin{solution}{exo:b2:mantle-melting:6}
(a) Forsterit: $60.08/(2\times40.30 + 60.08) = 60.08/140.68 =
\qty{42.71}{\percent}$ \ce{SiO2}; enstatit: $60.08/100.38 = \qty{59.85}{\percent}$.
(b) $0.975\times59.85 + 2.5 = \qty{60.85}{\percent}$. (c) Enstatit terletak di antara
forsterit dan cairan: fraksi forsterit adalah $(60.85 - 59.85)/(60.85 -
42.71) = 1.00/18.14 = \qty{5.5}{\percent}$, sama dengan nilai Bowen dan Andersen. (d) Lelehan
lebih kaya akan silika daripada ortopiroksen asalnya, dan residu
mendapat tambahan olivin.
\end{solution}

\begin{solution}{exo:b2:mantle-melting:7}
(a) $5.1P^2 - 112.9P + 214.3 = 0$: $P_0 = (112.9 - \sqrt{12746.4 -
4371.7})/10.2 = (112.9 - 91.51)/10.2 = \qty{2.10}{GPa}$, pada $2.097\times30.9 =
\qty{64.8}{km}$. (b) $5.1P^2 - 112.9P + 364.3 = 0$: $P_0 = (112.9 -
72.90)/10.2 = \qty{3.92}{GPa}$, pada \qty{121}{km}. (c) $T_{\mathrm s}'(2.10) = 132.9
- 21.4 = \qty{111.5}{K/GPa}$, sehingga $\Delta P_0 = 10/(111.5 - 20) = \qty{0.109}{GPa}$
per \qty{10}{\celsius}, atau \qty{3.4}{km}.
\end{solution}

\begin{solution}{exo:b2:mantle-melting:8}
(a) $T_{\mathrm s} = \qty{1331.1}{\celsius}$, $T_{\mathrm l} = 1475.0 + 160.0 -
12.8 = \qty{1622.2}{\celsius}$; $T = 1331.1 + 0.10\times291.1 =
\qty{1360.2}{\celsius}$, atau \qty{1633}{K}. (b) $(\partial T/\partial P)_F = 112.5
+ 0.10\times(67.2 - 112.5) = \qty{108.0}{K/GPa}$; $(\partial T/\partial F)_P =
T_{\mathrm l} - T_{\mathrm s} = \qty{291.1}{K}$. (c) $c_p/T = 0.612$; dengan $\alpha/\rho
= \qty{12.1}{J/(kg.K)}$ per gigapascal, pembilangnya $0.612\times108.0 - 12.1
= 54.0$ dan penyebutnya $300 + 0.612\times291.1 = 478.2$: $-\mathrm dF/\mathrm
dP = 0.113$, yaitu \qty{11.3}{\percent} per gigapascal. (d) $54.0/178.2 = 0.303$,
\qty{30}{\percent} per gigapascal: kalor laten memangkas produktivitas dengan
faktor 2.7.
\end{solution}

\begin{solution}{exo:b2:mantle-melting:9}
(a) $43\times1 = \qty{43}{\celsius}$; $43\times5^{0.75} = 43\times3.34 =
\qty{144}{\celsius}$; $43\times13^{0.75} = 43\times6.85 = \qty{294}{\celsius}$.
(b) $1213.5 - 294 = \qty{919}{\celsius}$. (c) \qty{125}{ppm} sama dengan
\qty{0.0125}{\percent}, sehingga $X = 0.0125/0.01 = \qty{1.25}{\percent}$ dan $\Delta T =
43\times1.25^{0.75} = 43\times1.18 = \qty{51}{\celsius}$.
\end{solution}

\begin{solution}{exo:b2:mantle-melting:10}
(a) Lebar dasar segitiga $2H$ dan mantel melintasinya dengan kecepatan $U$: $2UH$ per satuan panjang
punggungan dan per satuan waktu. Setiap sisi membentang setinggi $H$ dan dilintasi secara horizontal dengan kecepatan
$U$: masing-masing $UH$, seluruhnya $2UH$. Dengan sisi yang miring $\theta$, lebar dasarnya akan
$2H/\tan\theta$ dan aliran masuknya $2UH/\tan\theta$, yang sama dengan aliran keluar hanya jika
$\tan\theta = 1$: dalam model ini sudut itu ditentukan oleh kekekalan massa. (b) Dengan
produktivitas tetap, \omterm{def:b2:mantle-melting:fraction}{fraksi lelehan} di puncak adalah $F_{\max} = \Pi P_0 = \Pi\rho
g z_0$, dan $H = z_0$: $h_c = F_{\max}z_0/2 = \Pi\rho g z_0^2/2$. (c) Setiap tahun
terbentuk \qty{50}{mm} lempeng baru, yaitu $\qty{0.05}{m}\times\qty{1000}{m}\times\qty{7100}{m} =
\qty{3.6e5}{m^3}$ kerak per kilometer punggungan; dengan \omterm{def:b2:mantle-melting:fraction}{fraksi lelehan} rata-rata \qty{12}{\percent},
kerak itu berasal dari $\num{3.6e5}/0.12 = \qty{3.0e6}{m^3}$ mantel.
\end{solution}

\begin{solution}{exo:b2:mantle-melting:11}
(a) Setelah massa $t$ A meninggalkan satu satuan massa cairan, cairan mengandung $0.6 -
t$ A, 0.2 B, dan 0.2 C dalam massa $1 - t$: $\big((0.6 - t), 0.2,
0.2\big)/(1 - t)$, dan B/C tetap 1. (b) Pada kotektik, $(0.40 - 0.15s, 0.60 -
0.30s, 0.45s)$; B = C memberikan $0.60 - 0.30s = 0.45s$, $s = 0.8$: $L_1$ berkomposisi A 28, B
36, C 36. Maka $0.2/(1 - t) = 0.36$, $t = 0.44$: \qty{44}{\percent} sistem telah
mengkristal sebagai A. (c) Seluruh C berada dalam cairan: $f = 20/45 = 0.44$, jadi
\qty{44}{\percent} masih cair. (d) A: $0.60 - 0.444\times0.25 = 0.49$; B: $0.20 -
0.444\times0.30 = 0.07$; bersama cairannya, $49 + 7 + 44 = \qty{100}{\percent}$.
\end{solution}

\begin{solution}{exo:b2:mantle-melting:12}
(a) $X_{\mathrm{sat}}(1) = 12 + 1 = \qty{13}{\percent}$, $\Delta T =
\qty{294}{\celsius}$, solidus jenuh $1213.5 - 294 = \qty{919}{\celsius}$;
$X_{\mathrm{sat}}(2) = 12\times1.516 + 2 = \qty{20.2}{\percent}$, $\Delta T =
43\times20.2^{0.75} = 43\times9.53 = \qty{410}{\celsius}$, solidus jenuh $1331.1
- 410 = \qty{921}{\celsius}$. (b) Lelehan melarutkan lebih banyak air pada tekanan yang lebih tinggi,
dan penurunan yang lebih besar itu mengimbangi kenaikan solidus kering sebesar \qty{118}{\celsius}.
(c) Solidus harus turun sebesar $1331.1 - 1100 = \qty{231}{\celsius}$: $X =
(231/43)^{1/0.75} = 5.37^{1.333} = \qty{9.4}{\percent}$ air, kurang dari
\qty{20.2}{\percent} pada keadaan jenuh, sehingga baji itu dapat meleleh.
\end{solution}

\begin{solution}{pb:b2:mantle-melting:1}
\textbf{1.} $1350 + 273 = \qty{1623}{K}$.
\textbf{2.} $\num{4.0e-5}\times1623/(3300\times1000) = \qty{1.97e-8}{K/Pa} =
\qty{19.7}{K/GPa}$, dibulatkan menjadi \qty{20}{K/GPa}.
\textbf{3.} $\num{e9}/(3300\times9.81) = \qty{30.9}{km}$; $20/30.9 =
\qty{0.65}{K/km}$.
\textbf{4.} $1350 + 20P = 1085.7 + 132.9P - 5.1P^2$, yaitu $5.1P^2 - 112.9P +
264.3 = 0$: $P_0 = (112.9 - \sqrt{12746.4 - 5391.7})/10.2 = (112.9 - 85.76)/10.2 =
\qty{2.66}{GPa}$.
\textbf{5.} $2.66\times30.9 = \qty{82}{km}$; $T = 1350 + 20\times2.66 =
\qty{1403}{\celsius}$.
\textbf{6.} $1085.7 + 132.9\times2.66 - 5.1\times2.66^2 = 1085.7 + 353.5 - 36.1 =
\qty{1403}{\celsius}$.
\textbf{7.} $T_{\mathrm s} = 1085.7 + 199.4 - 11.5 = \qty{1273.6}{\celsius}$;
$T_{\mathrm l} = 1475.0 + 120.0 - 7.2 = \qty{1587.8}{\celsius}$; selisihnya
\qty{314.2}{K}.
\textbf{8.} $1273.6 + 0.10\times314.2 = \qty{1305.0}{\celsius}$, \qty{1578}{K}.
\textbf{9.} $T_{\mathrm s}' = 132.9 - 15.3 = \qty{117.6}{K/GPa}$, $T_{\mathrm l}' =
80.0 - 9.6 = \qty{70.4}{K/GPa}$: $(\partial T/\partial P)_F = 117.6 +
0.10\times(70.4 - 117.6) = \qty{112.9}{K/GPa}$; $(\partial T/\partial F)_P =
\qty{314.2}{K}$.
\textbf{10.} $c_p/T = 1000/1578 = 0.634$; pembilang $0.634\times112.9 - 12.1 = 71.6
- 12.1 = 59.5$; penyebut $300 + 0.634\times314.2 = 499.2$: $0.119$, atau
\qty{11.9}{\percent} per gigapascal.
\textbf{11.} Dengan fungsi hasil pencocokan $F = T'^{1.5}$, dengan $T'$ bagian jarak yang telah ditempuh dari solidus ke likuidus, suhu harus naik tajam terhadap $F$ untuk menghasilkan
persen pertama lelehan, sehingga $(\partial T/\partial F)_P$ besar di dekat solidus
dan pelelehan dimulai dengan lambat. $F_{\max} = 0.09\times2.66 = 0.24$, yaitu \qty{24}{\percent}.
\textbf{12.} $0.13/0.50 = 0.26$: klinopiroksen habis pada pelelehan \qty{26}{\percent}.
Pada \qty{24}{\percent}, residu di bawah sumbu masih mengandung sedikit klinopiroksen:
\omterm{def:b2:mantle-melting:peridotites}{lherzolit} yang miskin klinopiroksen, di ambang \omterm{def:b2:mantle-melting:peridotites}{harzburgit}.
\textbf{13.} $h_c = 0.24\times82/2 = \qty{9.8}{km}$.
\textbf{14.} $F_{\max}/2 = \qty{12}{\percent}$.
\textbf{15.} \qty{9.8}{km} adalah \qty{38}{\percent} lebih tebal daripada \qty{7.1}{km}, di atas
rentang kerak normal.
\textbf{16.} Produktivitas tetap memberikan terlalu banyak lelehan kepada mantel yang paling dalam;
model hasil pencocokan menghasilkan lelehan pertama secara perlahan, $F(z)$ tetap lebih rendah di dekat awal pelelehan, dan
$\int F\,\mathrm dz$ lebih kecil.
\textbf{17.} $F_{\max} = \Pi P_0$ dan $H = 30.9P_0$ (km): $h_c = \Pi P_0\times30.9P_0/2
= kP_0^2$ dengan $k = 0.09\times30.9/2 = \qty{1.39}{km/GPa^2}$.
\textbf{18.} $P_0 = \sqrt{7.1/1.39} = \sqrt{5.11} = \qty{2.26}{GPa}$.
\textbf{19.} $2.26\times30.9 = \qty{70}{km}$.
\textbf{20.} $T_{\mathrm s}(2.26) = 1085.7 + 300.4 - 26.0 = \qty{1360.0}{\celsius}$,
dan $T_p = 1360.0 - 20\times2.26 = 1360.0 - 45.2 = \qty{1315}{\celsius}$.
\textbf{21.} Untuk \qty{6.3}{km}: $P_0 = \sqrt{4.53} = \qty{2.13}{GPa}$, $T_{\mathrm s}
= \qty{1345.6}{\celsius}$, $T_p = \qty{1303}{\celsius}$. Untuk \qty{7.9}{km}: $P_0 =
\sqrt{5.68} = \qty{2.38}{GPa}$, $T_{\mathrm s} = \qty{1373.1}{\celsius}$, $T_p =
\qty{1326}{\celsius}$. Sebaran tebal kerak mengizinkan $1315 \pm
\qty{12}{\celsius}$.
\textbf{22.} $P_0 = \sqrt{20/1.39} = \sqrt{14.39} = \qty{3.79}{GPa}$; $T_{\mathrm s} =
1085.7 + 503.7 - 73.3 = \qty{1516.1}{\celsius}$; $T_p = 1516.1 - 75.8 =
\qty{1440}{\celsius}$, kelebihan sebesar $1440 - 1315 = \qty{125}{\celsius}$.
\textbf{23.} Keduanya memakai bukti yang berbeda. Lava merekam mantel yang paling panas,
di pusat plume; kerak merata-ratakan lelehan seluruh mantel yang
melewati wilayah pelelehan, dan sebagian besar mantel itu lebih dingin. Model sederhana itu juga menganggap
produktivitas tetap, padahal klinopiroksen habis pada \qty{26}{\percent}
sedangkan di sini $\Pi P_0 = 0.09\times3.79 = 0.34$: mantel panas yang sebenarnya meleleh lebih sedikit, dan memerlukan
suhu yang lebih tinggi untuk menghasilkan kerak yang sama.
\textbf{24.} $T_{\mathrm s}'(2.26) = 132.9 - 23.1 = \qty{109.8}{K/GPa}$; $\Delta P_0 =
10/(109.8 - 20) = \qty{0.111}{GPa}$; karena $h_c \propto P_0^2$, $\Delta h_c/h_c =
2\times0.111/2.26 = \qty{9.8}{\percent}$, sekitar sepersepuluh per \qty{10}{\celsius}.
\textbf{25.} Di punggungan yang sangat lambat, mantel naik begitu lambat sehingga kehilangan kalor secara
konduksi ke permukaan yang dingin; mantel meninggalkan adiabat, pelelehannya berhenti
lebih dalam, dan bagian atas segitiga, tempat $F$ paling besar, terpangkas: keraknya
lebih tipis.
\textbf{26.} Kerak punggungan normal memberikan \omterm{def:b2:mantle-melting:tp}{suhu potensial mantel}
\textbf{\qty[text-series-to-math=true]{1315}{\celsius}}, dengan ketidakpastian sekitar
\qty{12}{\celsius}: \qty{35}{\celsius} di bawah \qty{1350}{\celsius} dari fisika
mineral, yang ketidakpastiannya kira-kira sebesar itu juga, dan jauh di dalam rentang 1280 sampai
\qty{1400}{\celsius} dari lava. Tiga termometer yang saling independen sepakat pada mantel
bersuhu sekitar \qty{1300}{\celsius} sampai \qty{1350}{\celsius}.
\end{solution}
